

\section{引言}

随着电力、交通、能源和大型工业设施持续向规模化、复杂化方向发展，巡检对象的空间尺度、结构复杂程度以及作业环境的危险性不断增加。输电线路、桥梁、风电机组、光伏场站、隧道、压力钢管和大型厂房等设施具有共同的特点：一部分设备分布范围广，另一部分设备则处于高空、狭窄或人员难以长期到达的区域；同时，设备状态具有连续变化特征，巡检需要在保证作业安全的条件下获得具有空间对应关系和足够质量的观测数据。传统人工巡检仍是许多场景的重要手段，但其作业效率、可达性以及结果一致性受到人员条件和现场环境的制约。无人机（Unmanned Aerial Vehicle, UAV）能够以非接触方式进入高空、狭窄和大范围区域，并可搭载可见光、红外、激光雷达等传感器，因此逐渐成为基础设施和工业设备巡检的重要数据获取平台\cite{Nguyen2018PowerLine,Ahmed2020BridgeNDE,Zhang2022UBI,Panigati2025Bridge,Yahya2022PVReview,Shakhatreh2019UAV,nooralishahi2021drones,safie2025trip,nooralishahi2022access}。

从应用演进看，无人机在工业巡检中的角色已经发生变化。早期应用主要关注飞行平台获取影像和测量数据的能力，巡检人员仍需要根据既定航线完成数据采集，并在飞行结束后进行人工分析。随着摄影测量、三维重建和计算机视觉的发展，研究逐步将设备识别、缺陷检测以及三维建模纳入同一工作流程\cite{Colomina2014UAV,Nex2014UAV,Jeong2020BridgeReview,Lin2021BridgeAerialRobots}。桥梁巡检领域的系统性研究表明，现有无人机巡检已经覆盖图像采集、三维建模、损伤识别和结果分析等多个环节，但完整流程仍常需要人工参与\cite{Zhang2022UBI,Lin2021BridgeAerialRobots}。输电线路巡检和光伏场站巡检的研究也表现出类似特点：无人机能够降低大范围数据采集的时间与人员风险，但检测对象、成像条件和现场运行状态仍会影响最终的巡检结果\cite{Nguyen2018PowerLine,Ahmed2024Transmission,Yahya2022PVReview,Ramirez2022PVRTK}。

因此，复杂工业场景下的无人机自主巡检不能简单理解为将人工巡检中的相机安装到飞行平台上。巡检任务本身包含至少三个相互关联的环节。首先，无人机需要确定自身位姿并建立与周围环境的空间关系，以保证飞行安全和观测位置的准确性；其次，需要围绕具体设备获取能够支持状态判断的信息，包括设备外观、几何结构、温度场以及与历史状态相关的变化信息；最后，巡检结果还应作用于后续飞行决策，例如对疑似异常区域进行近距离复查，对未充分观测区域重新规划视点，或者根据任务优先级调整巡检顺序。由此，巡检中的感知、异常识别和任务规划并不是相互独立的功能模块，而是具有明确的数据传递和反馈关系\cite{Zhang2022UBI,Alqudsi2026GNSSDenied}。

复杂工业场景带来的困难主要表现在两个方面。一方面，环境本身具有明显的非理想特征。室内厂房、隧道和设备内部可能不存在稳定的全球导航卫星系统（Global Navigation Satellite System, GNSS）信号，钢结构、管线和设备表面还可能形成遮挡、重复纹理和狭长通道；粉尘、烟雾、低照度、强反光等因素又会使视觉或激光观测出现不同程度的退化\cite{Ozaslan2017Penstocks,Chang2023GNSSDenied,Zhao2024ResilientUAV,jarraya2025gnssdenied,boiteau2024uavlowvisibility}。另一方面，巡检对象具有较强的任务属性。无人机不是为了单纯获得一张完整的环境地图，而是需要在有限时间和能量条件下获取与设备状态判断直接相关的信息。对于输电线路，重点在于线路及部件的定位和缺陷识别；对于桥梁，重点在于结构表面病害和三维空间对应关系；对于光伏和风电设施，则需要结合可见光、热红外等数据识别不同类型的运行或结构异常\cite{Nguyen2018PowerLine,Jeong2020BridgeReview,Yahya2022PVReview,Liu2022WindRobots,Sun2022WindMonitoring}。

现有研究已经在上述各个环节形成了较为成熟的技术基础，但不同研究关注的问题仍有明显差异。导航研究通常以定位误差、建图精度和系统实时性为主要评价指标；缺陷检测研究则主要关注识别准确率、召回率和模型泛化能力；巡检规划研究更多考虑路径长度、任务时间、覆盖率和能耗。针对完整自主巡检过程，如何将不同层面的信息统一起来，仍是研究中的重要问题。近期关于无人机韧性自主定位、多传感器融合以及全流程桥梁巡检的研究均指出，长期自主运行不仅取决于单一算法性能，还受到传感器退化、计算资源、任务耦合和人工干预程度的共同影响\cite{Zhao2024ResilientUAV,Li2025UAVFusion,Zhang2022UBI,Alqudsi2026GNSSDenied,zhao2024uavresilient,zhao2025resilient}。

基于上述认识，本文按照巡检过程中的实际任务对相关研究进行组织第二章首先从工业场景与巡检任务出发，分析无人机在不同应用环境中承担的具体工作以及由此产生的信息需求；随后讨论视觉、激光和惯性等传感器在环境感知、自主定位和建图中的应用，并分析单一传感器在复杂环境中的适用边界；最后从数据、特征和状态估计三个层面梳理多模态信息融合方法，重点讨论不同模态的互补关系以及融合方法在工业巡检中的实际作用。第三章进一步讨论设备缺陷与异常检测，第四章转向巡检任务规划和多无人机协同，最终从感知、认知和决策之间的关系总结当前研究仍未解决的问题，并提出后续研究方向。

\section{复杂工业场景无人机环境感知与信息融合}

无人机巡检任务具有明显的场景依赖性。不同设施的空间尺度、结构形式、缺陷类型和运行条件不同，对无人机提出的要求也不同。与其将这些场景笼统归为室外、室内等几类，更有必要从巡检任务和感知条件出发分析其差异：在大范围设施中，需要解决覆盖和目标搜索问题；在受限空间中，需要解决GNSS不可用条件下的定位与避障问题；在视觉或激光容易退化的环境中，则需要利用不同传感器的互补信息维持稳定感知。基于这种划分，下面依次讨论场景与任务、环境感知以及多模态融合。

\subsection{复杂工业场景及巡检任务分类}

无人机在工业和基础设施巡检中的应用范围已经覆盖输电线路、桥梁、建筑外立面、隧道、压力钢管、风电机组和光伏场站等多个领域\cite{Nguyen2018PowerLine,Jeong2020BridgeReview,Falorca2020Facade,Ozaslan2017Penstocks}。这些应用虽然属于不同产业，但从无人机执行任务的角度可以归纳为三类典型场景。

第一类为大范围开放或半开放设施。输电线路、桥梁、风电场和光伏场站通常具有空间范围大、目标分散和巡检区域长等特点。无人机的任务首先是沿设施分布方向搜索并覆盖目标区域，在此基础上获取设备及关键部件的高分辨率图像或热红外信息。输电线路巡检研究将任务进一步分解为线路及部件定位、缺陷识别、植被侵限监测和灾害巡查等内容\cite{Nguyen2018PowerLine,Ahmed2024Transmission}。光伏场站则通常需要对大面积组件进行连续观测，并结合可见光、红外以及其他成像手段识别热异常、裂纹和组件缺陷\cite{Yahya2022PVReview,Ramirez2022PVRTK,oliveira2023pvairt,zefri2023pvlwir}。 在此类任务中，巡检效率并不仅由无人机飞行速度决定，还受到覆盖完整性、观测距离、相机视场以及成像质量的共同影响。已有光伏无人机巡检研究已经将飞行高度、视角、航点和检查路线作为联立优化变量，说明数据采集轨迹本身就是巡检任务的一部分\cite{Vergara2022PVAerial}。

第二类为受限或封闭工业空间。隧道、压力钢管、管廊以及大型工业设备内部通常存在狭长空间、遮挡和较弱外部定位条件，其中部分环境还存在明显的低照度或纹理不足问题。Özaslan等在压力钢管和隧道中的研究表明，这类环境不能依赖外部全球定位系统（Global Positioning System, GPS）或固定摄像头完成自主飞行，而需要利用机载惯性测量单元（Inertial Measurement Unit, IMU）、激光和相机进行状态估计和环境建图\cite{Ozaslan2017Penstocks,yu2021rcar,alsayed2021rs}。相关后续工作进一步针对广义圆柱形空间的局部建图和状态估计进行了研究\cite{Ozaslan2018Cylinders}。受限空间巡检的核心任务因此从大范围覆盖转向沿设备结构稳定运动、保持安全距离并持续获得可用于缺陷判断的观测数据。其任务约束更接近近距离作业，而不是普通环境中的全局路径导航。

第三类为具有明显感知退化特征的工业环境。这类场景可以出现在前述两种场景中，其主要特征不是空间尺度，而是传感器的观测条件不稳定。例如，低照度、强反光、粉尘和烟雾会影响视觉质量；结构重复或纹理缺失会降低特征匹配可靠性；障碍遮挡会减少激光点云的有效信息。对于此类环境，仅根据单一传感器在标准数据集上的性能选择感知方案并不充分。无人机韧性自主定位研究指出，大范围复杂环境和长周期任务会同时放大感知退化、状态估计误差和资源限制带来的影响\cite{Zhao2024ResilientUAV}。这意味着工业巡检需要关注的不只是平均定位精度，还包括感知系统在局部失效或数据质量下降时能否维持连续运行\cite{Zhao2024ResilientUAV,boiteau2024uavlowvisibility}。

从巡检任务看，不同场景虽然对象不同，但任务可以进一步归纳为三个层次。第一层是环境与平台感知，包括自身位姿估计、障碍物感知、可通行空间判断以及局部地图更新。这些任务直接服务于飞行安全和路径执行。第二层是巡检对象感知，包括设备搜索、部件识别、目标定位和观测区域确定，为后续缺陷识别建立空间索引。第三层是状态信息获取，包括可见光、热红外、三维几何以及历史巡检信息的联合利用，用于判断设备是否存在异常。桥梁、输电线路和光伏巡检的研究虽然对象差异明显，但均包含从目标定位到状态判断的连续过程\cite{Zhang2022UBI,Nguyen2018PowerLine,Yahya2022PVReview,nooralishahi2021drones}。

这样的任务划分也说明了环境感知和设备检测之间的联系。设备检测需要满足一定的成像距离、视角和分辨率，环境感知则决定无人机能否到达并稳定保持相应的观测状态。如果只评价环境感知模块的定位精度，而不考虑最终获得的设备观测质量，就难以反映系统的真实巡检能力。桥梁无人机巡检研究已经开始从单项图像检测转向完整的数据采集、三维建模、缺陷分析和报告生成流程\cite{Lin2021BridgeAerialRobots,Zhang2022UBI}。这类工作所体现出的系统化趋势，为后续讨论自主任务规划提供了基础。

从信息需求角度看，不同传感器承担的功能并不相同。可见光相机能够提供纹理、颜色和丰富的语义信息，适合设备及表观缺陷识别；激光雷达（Light Detection and Ranging, LiDAR）直接提供三维几何信息，适合定位、建图和空间结构感知；IMU能够提供高频角速度和加速度信息，在快速运动和短时视觉退化条件下具有补偿作用；GNSS主要提供全局位置约束，在开阔环境中可以减少长期漂移；红外热像则能够记录温度分布，对部分电气和热状态异常具有不可替代的信息价值\cite{Hall1997Fusion,Khaleghi2013Fusion,Yahya2022PVReview,Ramirez2022PVRTK}。因此，工业巡检中的感知需求并不是简单增加传感器数量，而是根据具体任务确定需要获得的几何、运动、视觉和状态信息。

\subsection{工业巡检环境感知}

环境感知首先需要解决无人机的定位问题。对于室外开放环境，GNSS能够提供稳定的全局位置参考，是工程应用中最直接的导航信息来源；但在室内、隧道、大型工业设施内部以及建筑密集区域，GNSS信号可能受到遮挡、衰减或多径影响，无法持续提供可靠约束\cite{Chang2023GNSSDenied,Zhao2024ResilientUAV,Alqudsi2026GNSSDenied,jarraya2025gnssdenied,zhao2024uavresilient}。因此，自主定位研究长期围绕视觉、激光和惯性等机载传感器展开。对于无人机巡检而言，这一变化具有直接的工程含义：当外部定位基础不可靠时，无人机需要利用周围环境本身建立相对位置和空间结构关系。

视觉是无人机环境感知中应用最早、最广泛的传感器之一。相机具有质量小、成本低和信息量大的特点，可以同时提供几何和语义线索。视觉里程计（Visual Odometry, VO）与同步定位与建图（Simultaneous Localization and Mapping, SLAM）通过连续图像估计相机运动并构建环境地图，为GNSS拒止环境下的自主飞行提供了基础\cite{Scaramuzza2011VO,Cadena2016SLAM}。ORB-SLAM建立了较为完整的特征点视觉SLAM框架，将跟踪、局部建图、重定位和回环检测纳入统一系统\cite{MurArtal2015ORBSLAM}；ORB-SLAM2进一步扩展到单目、双目和RGB-D相机，并在包括工业环境在内的多类数据集和飞行平台上验证了系统性能\cite{MurArtal2017ORBSLAM2}；ORB-SLAM3又将视觉、视觉--惯性和多地图模式整合到统一框架中，提高了系统在不同传感器配置下的适用性\cite{Campos2021ORBSLAM3}。这些研究奠定了视觉SLAM在无人机自主定位中的基础地位。

视觉方法的局限同样十分明确。其性能依赖于纹理分布、光照条件和相机运动状态，在低照度、强反光、快速运动或弱纹理区域，特征数量和匹配稳定性会降低。对于工业设施而言，这些情况并不罕见。例如压力钢管和隧道内部通常具有暗环境和长距离重复结构，单纯依赖视觉难以稳定完成定位\cite{Ozaslan2017Penstocks}。因此，研究逐步转向视觉与惯性信息的联合利用。视觉-惯性导航系统（Visual-Inertial Navigation System, VINS-Mono）通过预积分IMU测量和视觉特征的紧耦合优化实现单目视觉--惯性状态估计，并进一步支持无人机闭环自主飞行\cite{Qin2018VINS}。视觉与惯性的互补主要体现在时间尺度上：IMU能够提供高频运动约束，视觉则提供与环境结构相关的几何信息，两者结合能够减弱短时间视觉退化对状态估计的影响。

激光雷达提供的距离信息使其对光照变化不敏感，因此成为复杂工业环境中视觉感知的重要补充。激光里程计与建图（Lidar Odometry and Mapping, LOAM）通过将激光里程计与地图构建结合，实现了实时三维激光建图\cite{Zhang2014LOAM}。随着计算能力提高，激光--惯性融合逐渐成为无人机实时定位的重要技术路线。快速激光-惯性里程计（Fast LiDAR-Inertial Odometry, FAST-LIO2）采用紧耦合迭代滤波框架，直接利用原始点云与IMU测量进行状态估计和地图更新，减少了对手工几何特征提取的依赖\cite{Xu2022FASTLIO2,shan2020liosam,wu2025dalislam,kim2024adaptiveuavlio}；快速激光-惯性里程计（Faster LiDAR-Inertial Odometry, Faster-LIO）则利用并行稀疏体素等数据结构降低高频激光--惯性里程计的计算开销\cite{Bai2022FasterLIO}。这些方法对工业巡检的意义在于，复杂设备和建筑结构往往具有稳定的几何边界，即使视觉纹理较弱，也可以通过激光几何信息维持定位和地图更新。

对于同时具有视觉、激光和惯性信息的无人机平台，研究进一步发展到三传感器紧耦合。激光-视觉-惯性同步定位与建图（Lidar-Visual-Inertial SLAM, LVI-SAM）通过因子图将视觉--惯性和激光--惯性系统进行统一融合，并利用激光提供的深度信息改善视觉特征估计，同时使用视觉结果辅助激光匹配\cite{Shan2021LVISAM}。其优势不只是增加了一种传感器，而是允许不同模态在不同失效条件下互相提供约束。2025年的多传感器融合SLAM系统综述进一步将LiDAR--IMU、视觉--IMU、LiDAR--视觉以及LiDAR--IMU--视觉四类体系进行了系统比较，并指出单传感器SLAM在复杂环境中的稳定性受限于自身传感机理，多传感器融合的主要价值在于利用互补测量降低系统对单一模态的依赖\cite{Fan2025FusionSLAM,lin2022r3live,zheng2025fastlivo2,pfreundschuh2024coinlio,zhang2024lviofusion}。 对无人机而言，这种发展方向与工业巡检中感知条件不断变化的特点相吻合。

环境感知还需要处理动态障碍物和地图持续更新问题。传统SLAM通常假设环境在建图过程中基本静态，而工业厂区中的人员、车辆、吊装设备以及临时设施会导致局部场景变化。对于近距离巡检，无人机还需要持续保持与设备表面的安全距离，因此局部几何变化直接影响飞行轨迹。近年来的无人机自主导航综述已经将动态环境、长期漂移、故障检测以及机载资源约束列为GNSS拒止环境中的主要问题\cite{Chang2023GNSSDenied,Alqudsi2026GNSSDenied,wang2024dynamicsurvey}。这意味着环境感知的输出不应仅是一条离线轨迹，而应持续提供可供规划器使用的空间状态和可靠性信息。

从工程应用来看，工业巡检中的环境感知还具有明显的任务导向。对于压力钢管和隧道，主要需求是保持稳定的局部定位并沿结构表面完成覆盖；对于桥梁和建筑，则需要建立与结构表面对应的三维空间关系，以支持后续影像拼接和缺陷定位；对于输电线路和光伏场站，还需要在大范围飞行中保持全局位置一致性，以便将巡检结果映射回具体设备\cite{Ozaslan2017Penstocks,Zhang2022UBI,Ahmed2024Transmission,Yahya2022PVReview,bolourian2020lidarbridge,xu2023tunnelinspection}。因此，环境感知的评价指标不能完全脱离应用任务。仅在通用SLAM数据集上取得较低轨迹误差，并不能说明系统一定能够满足工业设备近距离观测的稳定性和重复性要求。

现有研究已经形成了从GNSS辅助导航、视觉SLAM、视觉--惯性、激光SLAM到多传感器紧耦合融合的技术链条。其发展动力并非算法形式本身发生变化，而是实际运行条件不断提出新的约束：从开放环境进入GNSS拒止区域后，需要解决自主定位；从纹理丰富环境进入低照度或弱纹理环境后，需要引入激光和惯性；当单一传感器的退化不可避免时，则需要利用多模态冗余提高系统连续性。当前问题在于，这种多模态感知仍较多围绕定位和建图展开，与设备检测所需要的目标可见性、观测质量以及异常识别需求衔接不足。因此，有必要进一步分析不同模态信息如何在统一任务下进行组织和融合。

\subsection{多模态信息融合}

多模态信息融合的研究基础在于不同传感器对同一对象提供具有差异性的观测。多传感器数据融合的经典研究已经指出，融合的目的不是简单增加数据量，而是在数据存在冗余、互补和不确定性的情况下获得更完整的状态描述\cite{Hall1997Fusion,Khaleghi2013Fusion}。对于无人机工业巡检，这种互补性尤其明显：视觉能够提供纹理和语义，激光能够提供几何结构，IMU提供高频运动信息，GNSS提供全局位置基准，红外则提供温度分布。不同模态对应不同物理量，因此适用的融合方式和产生的收益也存在差异。

按照信息处理所在位置，多传感器融合通常可分为数据级、特征级和决策级三类。数据级融合直接处理原始测量，信息保留最充分，但对时间同步、空间标定和传感器精度要求较高；特征级融合先分别提取各模态特征，再进行联合表达或状态估计；决策级融合则在各传感器完成局部识别或估计后，对多个结果进行组合\cite{Hall1997Fusion,Khaleghi2013Fusion}。 在无人机自主定位中，这些层次又表现为松耦合和紧耦合两种工程架构：松耦合方案通常先分别得到视觉、激光和惯性估计结果，再进行融合；紧耦合方案则将多源测量共同纳入统一的滤波、优化或因子图模型。随着计算资源和算法工具的发展，后者逐渐成为高性能自主定位系统的重要形式\cite{Fan2025FusionSLAM,Li2025UAVFusion}。

视觉--惯性融合是无人机多模态感知中较成熟的应用。VINS-Mono将视觉特征和预积分IMU信息放入统一优化框架，使相机与惯性测量在同一状态估计问题中相互约束\cite{Qin2018VINS}。 这类方法解决的并不是单纯的传感器精度问题，而是不同传感器时间尺度和误差累积特性的互补：视觉能够从环境中获得空间结构约束，但帧率有限且容易受成像质量影响；IMU更新频率高，但积分误差随时间累积。将两者联合后，系统可以在短时视觉质量下降时利用惯性维持状态连续性，再利用后续视觉观测抑制惯性漂移。

LiDAR--惯性融合进一步体现了几何测量与运动测量之间的互补关系。FAST-LIO2通过紧耦合迭代卡尔曼滤波直接处理原始点云和IMU测量，在保证实时性的同时提高了复杂环境中的状态估计能力\cite{Xu2022FASTLIO2}。Faster-LIO则进一步优化稀疏体素数据结构，提高了激光--惯性估计在高频运行条件下的计算效率\cite{Bai2022FasterLIO}。 对工业巡检而言，激光--惯性方案尤其适用于结构几何信息丰富而视觉纹理不足的环境，如隧道、管廊和大型设备内部。其不足也比较直接：激光雷达的成本、体积、能耗和计算需求高于普通相机，因此需要根据巡检任务确定是否值得采用。

LiDAR、视觉和IMU的联合融合进一步提高了复杂环境中的感知冗余。LVI-SAM通过两个相互关联的子系统进行视觉--惯性和激光--惯性状态估计，并在回环和地图构建过程中进行信息交互\cite{Shan2021LVISAM}。这种架构能够在部分传感器性能下降时利用其他模态继续提供约束。近期三传感器融合SLAM综述也表明，LiDAR、视觉和IMU已经成为多模态SLAM中最常见的传感器组合之一\cite{Fan2025FusionSLAM}。 对工业巡检而言，这类系统的实际价值还在于可以将三维结构与二维纹理联系起来：激光提供设备表面几何位置，视觉提供表面纹理和语义，从而为后续缺陷定位建立共同空间参考。

多模态融合的对象也逐渐从定位传感器扩展到巡检状态信息。光伏场站研究表明，热红外影像能够识别部分组件异常，而可见光图像则更适合表面和几何缺陷分析；两类信息具有明显的互补性\cite{Yahya2022PVReview,Ramirez2022PVRTK,oliveira2023pvairt,zefri2023pvlwir,li2024multispectraldefect}。 对大型其他大型能源设施的自动化状态评估同样呈现多源信息互补的特点，实际应用中常需要结合不同成像与测量手段获得更完整的状态信息。 这说明，在工业巡检中，多模态融合不应只理解为用于估计无人机姿态的传感器组合，还包括设备本身的多源状态观测。对于某些异常，仅依靠RGB图像并不能得到足够的信息，而引入温度、深度或三维结构后，异常的可观测性可能得到改善。

近年来，语义信息开始进入环境感知和地图构建过程。对于工业设备，几何地图只能说明设备的空间形态，而巡检系统还需要知道某一区域属于什么设备、什么部件以及此前是否出现过异常。数字孪生和智能设施管理研究强调将物理对象、空间模型和运行状态建立持续关联\cite{Tao2019DigitalTwin}。 这一思路与无人机巡检中的多模态感知能够形成对应关系：无人机获得的图像、点云和热像并不是相互独立的文件，而可以共同关联到具体设备对象及其空间位置。这样形成的语义化巡检信息，能够直接服务于后续缺陷检测和任务规划。

多模态融合还需要处理数据质量变化。工业现场的传感器并不总处于理想工作状态，例如相机可能出现过曝、欠曝或运动模糊，激光点云可能因为遮挡而稀疏，GNSS可能发生间歇性失锁。固定权重的融合方案通常默认各模态的可靠性相对稳定，这一假设在实验室条件下较容易满足，而在复杂工业环境中并不成立。无人机韧性自主定位研究已经将感知退化、鲁棒估计和多源冗余信息融合放在同一框架中讨论\cite{Zhao2024ResilientUAV,zhao2025resilient,jarraya2025gnssdenied}。 2026年的GNSS拒止无人机导航综述则进一步将传感器退化、动态场景、长期漂移和机载资源约束作为影响自主运行的共性问题\cite{Alqudsi2026GNSSDenied}。 因而，工业巡检中的融合策略需要考虑模态可靠性随时间和空间变化的情况。

从目前研究情况看，多模态融合已经从简单的测量组合发展到紧耦合状态估计、跨模态特征表达以及语义信息关联，但仍存在两个方面的不足。一方面，多数方法的主要目标仍是提高定位、建图或单项检测性能，融合层输出与巡检任务之间缺乏明确接口；另一方面，不同模态的可靠性通常由算法结构隐式处理，尚缺少与具体巡检任务相结合的质量评估和信息选择机制。对于实际巡检任务，真正需要判断的不是某一种传感器是否更先进，而是在当前位置和当前任务下，什么信息能够有效降低设备状态判断的不确定性。由此，多模态融合的研究重点有必要进一步从传感器层面的互补转向任务层面的信息价值。

本章的讨论表明，复杂工业场景的环境感知已经具备较成熟的技术基础，但其研究重点正逐渐从单纯定位扩展到面向巡检任务的综合感知。场景变化决定传感器需求，传感器互补推动融合方法发展，而最终的评价标准仍应落到巡检任务能否可靠完成。下一章将在此基础上讨论设备缺陷与异常检测，重点分析无人机获得的视觉、红外及其他多模态信息如何转化为设备状态判断，以及缺陷样本不足和未知异常对现有检测方法带来的限制。

\section{工业设备缺陷与异常检测}

第二章讨论了无人机如何获得自身状态、环境结构以及巡检对象的多源信息。对于工业巡检而言，感知信息的价值最终需要体现在设备状态判断上。无人机搭载的可见光、红外和三维传感器能够持续获得设备表面和局部空间信息，但原始数据本身并不能直接回答设备是否存在缺陷以及缺陷处于何种状态。因此，缺陷与异常检测构成环境感知向设备状态认知过渡的关键环节。

工业设备检测与一般自然场景目标识别存在明显差异。巡检对象通常具有相对固定的结构和类别，真正需要识别的是局部、细微且低频出现的状态变化。例如，桥梁表面的裂缝和剥落、输电线路部件的缺陷以及光伏组件的局部热异常，其面积可能只占整幅图像很小一部分，同时又容易受到观测距离、视角、光照和背景结构的影响\cite{Nguyen2018PowerLine,Jeong2020BridgeReview,Zhang2022UBI,Yahya2022PVReview,hua2025aerialsmall,zhang2025lowaltitudeuav,zhu2021visdrone}。 从数据角度看，正常样本通常远多于缺陷样本，且实际缺陷类型难以事先穷尽。工业异常检测研究正是在这种数据条件下逐渐形成了区别于传统监督分类和目标检测的方法体系\cite{Diers2023IJCVIQC,Liu2024IADSurvey,Li2025IADSurvey}。

就巡检任务而言，缺陷识别可以分为图像级异常判断、目标级缺陷检测和像素级缺陷分割三个层次。前者主要判断当前观测区域是否存在异常，适用于大范围快速筛查；目标检测进一步给出缺陷所在位置及类别，适合确定需要关注的具体部件；像素级分割则提供更精细的缺陷轮廓，为裂缝长度、腐蚀面积等定量评价提供条件\cite{Diers2023IJCVIQC,Taox2021Survey,shukla2025iadsurvey,li2025iadreview}。 三者并不是互相替代的关系，而是对应由粗到细的不同巡检需求。对于无人机巡检系统，更合理的工作方式通常是先利用低成本的图像级判断进行筛查，再对疑似区域进行目标定位和精细分割，从而减少无人机在非重点区域的无效观测。

\subsection{基于视觉与多模态信息的工业缺陷检测}

工业缺陷检测的发展首先得益于视觉数据的规模化获取。传统机器视觉通常依靠边缘、纹理、颜色和几何规则完成缺陷判定，在目标位置和成像条件相对稳定的生产线上仍具有应用价值，但对于复杂背景、光照变化和缺陷形态变化，其适应能力受到特征设计的限制\cite{Diers2023IJCVIQC,Li2024SurfaceDefectReview}。 深度学习通过从数据中自动学习特征表示，逐渐成为工业视觉检测的重要技术路线。相关综述表明，近年来的研究已经覆盖分类、检测和分割等多个任务，并进一步向少样本、不平衡数据和实时检测方向扩展\cite{Taox2021Survey,Diers2023IJCVIQC,Li2024SurfaceDefectReview}。

在有足够缺陷标注的条件下，监督式目标检测仍然是工业检测的主要方法之一。基于区域卷积神经网络（Faster Region-based Convolutional Neural Network, Faster R-CNN）采用区域提议网络实现候选区域生成与检测器联合训练，在复杂背景下具有较好的定位能力\cite{Ren2017FasterRCNN}。 YOLO 系列则以单阶段检测为代表，在推理速度方面更适合实时视觉应用\cite{Redmon2016YOLO,Redmon2018YOLOv3,Bochkovskiy2020YOLOv4,tan2024sccmdet,yi2025uavtransmission}。 单发多框检测器（Single Shot MultiBox Detector, SSD）同样通过多尺度特征进行单阶段检测，为不同尺寸目标提供统一的检测框架\cite{Liu2016SSD}。 对于工业巡检而言，检测器的选择并不能脱离无人机平台的运行条件：远距离飞行有利于提高覆盖效率，却会使缺陷在图像中占据更少的像素；近距离观测能够提高分辨率，但同时增加飞行时间、碰撞风险和能源消耗。因此，检测模型不仅要追求精度，还要与无人机观测距离和任务时间共同考虑。

工业缺陷具有明显的尺度差异和类别不平衡特征。裂纹等细长结构在像素层面往往具有较小宽度，而较大的剥落区域又可能覆盖复杂背景。Focal Loss 通过降低易分类样本对损失函数的贡献，使训练过程更加关注难分类样本，为类别不平衡问题提供了具有代表性的解决思路\cite{Lin2020Focal}。 多尺度特征融合则成为工业检测中的常见策略，其直接原因在于不同飞行距离、不同镜头焦距和不同设备尺寸会造成目标尺度变化。对于无人机巡检，尺度问题还与航迹相关：同一设备在不同航段采集的图像可能对应完全不同的目标尺寸，因此检测模型的尺度鲁棒性不能只通过固定数据集上的平均指标进行判断。

随着 Transformer 进入视觉领域，工业检测又出现了新的网络架构。检测Transformer（DEtection TRansformer, DETR）将目标检测转化为集合预测问题，减少了传统检测流程中部分手工设计的组件\cite{Carion2020DETR}。 Deformable DETR 针对高分辨率和多尺度目标引入可变形注意力，改善了检测计算效率和小目标建模能力\cite{Zhu2021DeformableDETR}。 视觉Transformer（Vision Transformer, ViT）和 Swin Transformer 则分别提供了全局和层次化的视觉特征表示\cite{Dosovitskiy2021ViT,Liu2021Swin}。 这些研究对工业缺陷检测的启示主要在于，缺陷判断不仅依赖局部纹理，也需要结合设备整体结构和局部区域之间的关系。但对于无人机平台而言，模型复杂度仍然是实际部署的重要约束。高分辨率输入、较大的模型参数量和持续推理都会增加机载计算功耗，因此工业巡检中常见的精度提升不能简单等同于系统能力提升。

围绕这一矛盾，轻量化检测方法成为工程应用的重要方向。MobileNet 通过深度可分离卷积降低计算量，EfficientNet 通过网络深度、宽度和输入分辨率的联合缩放实现精度与计算量之间的折中\cite{Howard2017MobileNet,Tan2019EfficientNet}。 对无人机巡检而言，轻量化模型的意义不仅是提高帧率，还包括在有限电池容量下延长有效巡检时间。尤其对于大范围输电线路和光伏场站巡检，若检测算法需要长时间运行，单位图像计算成本会直接转化为平台能源开销。因此，模型选择应放在“检测精度—图像分辨率—推理时间—能源消耗”的整体关系中考察。

在典型基础设施应用中，监督式深度学习已经表现出较好的缺陷识别能力。Yeum 和 Dyke 利用视觉方法研究桥梁裂缝自动检测，说明传统机器视觉和学习方法可以在复杂结构表面识别裂缝\cite{Yeum2015Vision,xiao2024bridgeroicrack}。 Cha 等采用卷积神经网络开展混凝土裂缝识别，进一步提高了自动特征提取能力\cite{Cha2017Deep,wang2024steelcorrosion}。 Liu 等将无人机影像与三维场景重建结合，用于桥墩裂缝评价，使缺陷检测结果能够与三维空间位置建立联系\cite{Liu2020BridgeCrack,maboudi2025bridgedeformation}。 在输电线路场景中，Tao 等利用卷积神经网络开展绝缘子缺陷检测，表明深度视觉方法已经能够处理具有明确结构和缺陷类型的电力设备巡检任务\cite{Tao2020Insulator,hang2024wtbcrack}。 这些研究共同表明，针对已知设备和已知缺陷，监督式视觉检测已经形成相对成熟的技术基础。

但是，监督式检测的适用性高度依赖训练数据。工业设备中的缺陷通常具有长尾分布，严重缺陷出现次数少，而正常状态样本数量远大于缺陷样本；同时，不同设备型号、拍摄距离和运行环境会导致图像分布发生变化\cite{Li2024Imbalance,Diers2023IJCVIQC}。 对无人机而言，这种分布变化更为突出，因为拍摄条件不仅受设备影响，还受飞行高度、姿态和视角变化影响。一个在固定视角数据集上训练得到的检测器，未必能够直接适应实际巡检中的连续视角变化。因此，监督检测虽然解决了“已知缺陷如何识别”的问题，但并没有解决“新设备、新场景以及未知缺陷如何处理”的问题。

多模态信息为这一问题提供了另一种思路。第二章已经指出，可见光、红外和三维信息具有不同的物理含义。在缺陷检测中，这种互补性可以进一步转化为状态判据。例如，可见光适合描述裂纹、腐蚀和表面剥落等外观变化，红外能够反映温度分布，对光伏组件热斑及部分电气设备异常具有较高敏感性\cite{Jadin2012Infrared,Usamentiaga2014Infrared,Ramirez2022PVRTK,santos2024uvthermopl,liu2025infraredcsee,ertunc2024thermalfault,li2024multispectraldefect}。 三维信息则能够描述局部几何变化，对于凹陷、变形、缺失和装配偏差等异常具有优势。2025 年关于无监督工业异常检测的综述已经将红绿蓝（Red-Green-Blue, RGB）、3D 及 RGB--3D 多模态方法分开讨论，说明工业异常认知正在由单一二维外观信息向外观与几何联合建模扩展\cite{Li2025RGB3D}。 

对于无人机巡检，多模态检测的价值还体现在观测条件发生变化时的补偿作用。视觉图像受到光照影响较大时，红外信息可能仍能提供状态差异；单纯依赖二维图像难以确定缺陷是否对应真实结构变化时，三维信息可以提供额外的几何约束。因此，多模态并不意味着将所有数据简单拼接，而应根据具体缺陷机理确定需要融合的观测信息。其难点在于不同传感器的成像时间、空间分辨率和坐标体系并不一致，同时不同模态对异常的敏感程度存在差异。对于无人机边缘平台，还需要额外考虑数据同步、通信带宽和算力等约束。由此，工业缺陷检测的研究重点逐渐从单一检测器性能，转向多源信息条件下的可靠状态判断。

\subsection{小样本、无监督异常检测}

真实工业环境中，异常数据难以大量获得是监督式缺陷检测面临的基本问题。设备正常运行时能够产生大量正常样本，而真实故障或缺陷具有低发生率，部分高风险设备又不适合通过人为制造缺陷获得训练数据。缺陷样本还往往具有很强的类别和设备依赖性，新设备可能出现训练阶段没有覆盖的异常形式。因此，工业异常检测逐渐形成了以正常样本建模为基础的研究范式，即利用正常数据学习设备或产品的正常分布，再根据测试样本偏离正常模式的程度判断是否异常\cite{Pang2021DeepAD,Diers2023IJCVIQC,Liu2024IADSurvey,li2025iadreview,shukla2025iadsurvey}。

异常检测与传统分类的一个重要区别在于，异常通常并不具有预先定义的完整类别集合。对于分类问题，训练目标是把样本映射到有限类别；而异常检测更关心测试样本是否偏离正常状态。由此，异常检测方法的核心任务从“学习不同缺陷之间的分类边界”转变为“建立可靠的正常性表示”。这一区别决定了其方法体系更多围绕重建、特征分布建模、知识蒸馏和自监督表示学习展开\cite{Pang2021DeepAD,Diers2023IJCVIQC,Li2025IADSurvey}。

早期无监督方法主要采用重建模型。变分自编码器通过学习潜在变量分布建立正常样本的生成模型，测试时利用重建误差衡量样本与正常模式的偏离程度\cite{Kingma2014VAE}。 生成对抗网络则利用生成器和判别器学习数据分布，为无监督异常检测提供了另一种生成式路线\cite{Goodfellow2014GAN}。 这类方法的共同特点是训练阶段只需要正常样本，但在工业应用中容易出现一个问题：模型具有较强的重建能力时，一些缺陷也可能被较好地重建，从而导致异常与正常区域之间的重建差异减小。由此，后续研究逐渐从图像重建转向更加稳定的特征空间建模。

深度单类分类是这一转变的重要方向。深度单类支持向量数据描述（Deep Support Vector Data Description, Deep SVDD）将正常样本映射到特征空间中的紧致区域，通过样本与正常中心之间的距离度量异常程度\cite{Ruff2018DeepSVDD}。 该方法从正常样本分布角度描述异常，而不是要求模型复原每一个像素，因此为工业异常检测提供了更加直接的正常性建模方式。随着 ImageNet 等大规模数据集预训练视觉特征的广泛应用，研究者进一步发现，具有较好语义表征能力的预训练网络可以直接为工业正常模式提供有用的特征空间。PANDA 通过适应预训练特征进行异常检测和分割，说明在缺陷标注不足的情况下，合理利用外部预训练知识能够减少从零学习的需求\cite{Reiss2021PANDA}。

MVTec异常检测（MVTec Anomaly Detection, MVTec AD）数据集的提出推动了这一方向的系统发展。该数据集包含多个工业对象和纹理类别，训练阶段主要使用无缺陷样本，测试阶段提供多种真实缺陷及像素级标注，为不同方法的检测和定位性能提供了统一基准\cite{Bergmann2019MVTec,Wang2024RealIAD}。 此后，PatchCore 等方法利用正常样本的局部特征建立记忆库，并通过测试样本与正常特征之间的相似度进行异常判断\cite{Roth2022PatchCore}。 PaDiM 则对局部深度特征建立概率分布，并通过统计距离进行异常定位\cite{Defard2021PaDiM}。 这些方法在工业基准上取得了较好的性能，反映出“预训练特征+正常样本建模”已经成为无监督工业异常检测的重要技术路线。

知识蒸馏和特征分布建模进一步降低了异常检测的推理成本。Deng 和 Li 提出的逆蒸馏方法以教师网络产生的单类表示为基础，通过学生解码器恢复多尺度特征，使异常区域表现为教师和学生表示之间的差异\cite{Deng2022RD}。 FastFlow 则利用二维归一化流对深度特征分布进行建模，避免显式构建复杂的非参数分布\cite{Yu2022FastFlow}。 EfficientAD 将教师--学生结构与轻量化设计结合，在强调工业异常定位能力的同时重点优化推理延迟\cite{Batzner2024EfficientAD,liu2023simplenet,coscia2025onen}。 对无人机巡检而言，实时性尤其重要，因为异常结果如果只能离线计算，就难以影响当前飞行状态；只有当检测延迟足够低时，异常信息才有可能反馈给视点选择和路径规划模块。

在样本进一步减少的情况下，自监督方法开始发挥作用。CutPaste 仅使用正常样本，通过裁剪和粘贴构造伪异常来训练异常判别表示，并在工业异常数据集上验证了这种方法对未知缺陷的识别能力\cite{Li2021CutPaste,zhu2024inctrl}。 DRAEM 则利用合成异常进行重建和判别联合学习，以提高异常区域定位能力\cite{Zavrtanik2021DRAEM,li2024musc}。 这类方法降低了真实缺陷数据的需求，但同时引入了新的假设，即构造出的伪异常应当与真实异常在局部结构或纹理层面具有一定相关性。当真实工业缺陷的物理机理与合成扰动存在明显差异时，模型的泛化能力仍然可能受到限制。

随着视觉基础模型的发展，少样本和零样本异常识别开始受到关注。对比语言-图像预训练（Contrastive Language-Image Pre-training, CLIP）通过大规模图像--文本对比学习获得跨模态视觉表征，为利用语言描述辅助视觉识别提供了通用基础\cite{Radford2021CLIP,zhou2024anomalyclip}。 WinCLIP 进一步将视觉语言模型引入工业异常检测，研究零样本和少正常样本条件下的异常分类与分割问题，在 MVTec AD 和 VisA 等数据集上取得了较好的结果\cite{Jeong2023WinCLIP,ma2025aaclip}。 这类方法的意义在于降低对每一种设备单独训练检测器的依赖，但工业异常本身往往具有细粒度、局部化和设备特定的特点，通用视觉语言知识与具体设备物理状态之间仍存在差距。因此，基础模型更适合作为异常认知体系的一部分，而不能直接替代设备相关的状态建模。

另一个值得注意的变化是，工业异常检测开始从二维 RGB 数据扩展到 3D 与多模态数据。MVTec 3D-AD 等数据集使模型能够利用表面深度和几何结构判断异常，而近期研究进一步系统比较了 RGB、3D 以及 RGB--3D 多模态异常检测的方法\cite{Bergmann2021MVTec3D,Li2025RGB3D,li2025mulsenad}。 这种方向与无人机巡检具有直接联系：对于从单张图像难以判断的变形、缺失和结构偏差，三维信息可以补充二维外观；对于热异常，则需要进一步引入红外等物理信息。由此，异常检测逐渐从单一图像上的统计偏离判断，向多源状态一致性判断发展。

但是，公开数据集上的高精度并不意味着算法已经能够适应真实工业现场。2024 年针对 AutoVI 数据集的研究指出，许多无监督异常检测方法主要在公开数据集上开发和评估，而这些数据集与真实工业生产条件之间存在明显差距，在更加真实的工业采集条件下，现有方法仍存在较大提升空间\cite{Carvalho2024AutoVI,shukla2025iadsurvey}。 2023 年工业图像质量控制综述同样指出，现有基准和评价指标仍存在不足，需要更多符合实际生产条件的数据集进行验证\cite{Diers2023IJCVIQC}。 对于无人机巡检，这一问题更加突出，因为设备图像同时受到飞行姿态、相机距离、背景变化、动态遮挡和环境光照等因素影响，公开固定视角数据集难以完整反映实际巡检过程。

因此，小样本和无监督异常检测虽然降低了对缺陷标注的依赖，但其真正的工程难点已经从“缺少异常标签”进一步转向“正常分布本身具有变化”。当设备型号、环境条件或采集方式变化时，正常样本的外观也会发生变化；如果模型不能区分正常工况变化与真实异常，就容易产生误报。另一方面，对于真正未知的异常，仅输出一个异常分数仍不足以支撑巡检决策。系统还需要回答异常位于何处、是否值得复查、需要采用何种传感器以及复查后结论是否发生变化。因而，面向无人机自主巡检的异常检测需要进一步考虑跨场景泛化、多模态状态融合以及检测结果与巡检行为之间的反馈关系。

综上，工业设备缺陷与异常检测已经形成从监督式视觉检测到无监督异常检测，再到基础模型与多模态异常认知的连续发展过程。监督式方法在设备类型和缺陷类别较明确时具有较强的检测能力，但依赖较多的标注数据；无监督和自监督方法降低了异常样本需求，并能够在一定程度上识别训练阶段未见过的缺陷；基础模型和多模态方法则进一步拓展了少样本、零样本和跨模态状态认知的能力\cite{Li2025IADSurvey,Li2025RGB3D,Mao2025IADSurvey,zhang2026continualad,zhu2024inctrl}。 当前仍需要解决的核心问题并不是单一模型的检测精度，而是异常判断的可靠性以及异常信息在自主巡检过程中的作用。对于无人机而言，只有当检测结果能够与环境位置、设备对象及后续观测行为建立联系，缺陷检测才真正成为自主巡检闭环中的状态认知环节。下一章将进一步讨论无人机如何利用这些感知和认知结果进行巡检任务规划、动态路径调整以及多无人机协同。

\section{无人机自主巡检规划与多无人机协同}

第二章和第三章分别讨论了复杂工业场景中的环境感知与信息融合，以及基于视觉和多模态信息的设备缺陷与异常检测。对于实际巡检任务而言，获得环境地图和设备状态判断并不是任务终点。无人机还需要确定巡检对象的访问顺序、观测位置和飞行轨迹，并在环境发生变化或检测结果不确定时调整原有计划。因此，巡检规划研究的重点已经由单纯寻找可行路径，逐渐转向同时考虑目标覆盖、观测质量、飞行安全、能源约束以及实时信息变化的任务级决策问题。

传统路径规划主要解决从起点到目标点的可行运动问题。采样型方法能够处理复杂障碍和高维运动约束，图搜索和组合优化方法便于显式加入距离、时间以及安全约束\cite{LaValle2006,Karaman2011,Kingston2018,mao2023powerinspection,bolourian2020lidarbridge,zhao2024structuralinspection,wang2025uavsurface}。但工业巡检与一般点到点导航存在明显差异：无人机需要对设备表面或特定空间区域进行有质量要求的观测，因此“是否完成有效观测”比“是否抵达目标位置”更能反映任务完成程度\cite{Galceran2013,Cabreira2019,Zeng2020,wang2025uavsurface,jung2024bridgepose,gazani2023bagofviews,bolourian2020lidarbridge}。当巡检范围进一步扩大，单架无人机又会受到续航和任务时间的限制，多无人机协同便成为提高覆盖效率的自然延伸\cite{Gerkey2004,Chakraa2023,Skaltsis2023,ivic2023multiuavinspection,gilcastilla2024hetero,jeong2023makespan,obrien2023dynamictask}。

\subsection{自主巡检任务规划与动态路径重规划}

对于大范围开放或半开放设施，工业巡检通常具有明显的覆盖属性。光伏场站、建筑外立面和部分大型设施需要对连续表面或区域进行系统观测，而桥梁、输电杆塔和风电设备则通常需要围绕特定部件布置观测位置。覆盖路径规划研究的基本目标是在保证区域覆盖的同时减少重复路径和任务代价\cite{Choset2001,Galceran2013,zhao2024structuralinspection}。Cabreira等系统总结了UAV覆盖路径规划研究，将覆盖率、路径长度、能耗和避碰等作为典型评价指标\cite{Cabreira2019}。对工业巡检而言，这些指标还需要与成像距离、视场范围和目标可见性结合，否则几何上的“覆盖”可能并不能转化为有效的巡检数据。

随着设施三维模型和数字化资料的使用，规划研究开始由区域覆盖转向面向目标表面的视点规划。视点规划首先根据目标表面、相机视场和距离约束生成候选观测位置，再确定这些位置的访问顺序和连接轨迹\cite{Zeng2020,wang2025uavsurface,jung2024bridgepose,gazani2023bagofviews,bolourian2020lidarbridge}。这种方法特别适用于桥梁、建筑外立面以及具有规则几何结构的大型设备。近期桥梁无人机巡检研究已经将三维模型、候选视点、碰撞风险、路径长度以及机动性约束同时纳入覆盖路径规划，说明巡检路径的评价标准正在从单纯的空间覆盖转向观测效果和飞行代价的综合权衡\cite{Panigati2025}。

三维模型还可以在规划开始阶段提供较强的先验信息。对于拥有建筑信息模型（Building Information Modeling, BIM）或其他三维模型的建筑和工业设施，候选观测点可以直接从目标表面生成，再结合相机和无人机约束形成初始巡检轨迹。这样能够降低现场搜索带来的计算负担，并使巡检路径与设备结构保持一致。但模型通常反映的是某一时刻的设施状态，临时障碍、人员和车辆运动以及设备状态变化难以提前写入模型。因此，离线模型驱动的规划更适合作为初始方案，实际执行仍需要在线调整\cite{ortega2020aiaa}。

当环境信息不完整，或者巡检系统希望主动获得更多有价值的信息时，下一最佳视点（Next-Best-View, NBV）方法成为重要研究方向。NBV根据当前地图中的未知区域、重建误差或信息增益选择下一观测位置，从而把规划目标从“到达某个点”转化为“获得具有较高价值的信息”\cite{Bircher2016,Zeng2020,gazani2023bagofviews,ruckin2022adaptiveipp}。 对工业巡检而言，这种机制尤其适合处理遮挡和异常复查。例如远距离图像只能判断某处存在疑似缺陷时，系统可以根据当前几何和视野条件重新选择观察位置，而不是继续沿预定航线飞行。

由此，巡检规划与第三章中的异常检测开始产生直接联系。检测模块给出的异常位置、异常置信度和结果不确定性都可以改变后续任务。当检测结果足够明确时，系统可以继续执行原任务；当置信度较低或设备风险较高时，则可能需要增加观测次数、缩短观测距离或者改变观察方向。换言之，巡检任务不再是任务开始时一次性给定的静态集合，而可能随着感知和检测结果不断发生变化。这种由感知结果产生新任务、再由新任务触发重新规划的机制，是自主巡检区别于固定航线作业的重要特征。

动态路径重规划主要用于处理环境和任务状态的持续变化。滚动时域方法并不一次性优化整条航迹，而是在每个时刻利用最新环境信息重新计算有限预测范围内的轨迹，因此可以在一定程度上平衡全局规划和在线计算\cite{Papaioannou2022,petit2024moar,ruckin2022adaptiveipp}。当现场出现临时障碍或者某些观测位置无法到达时，系统可以保留尚未完成的任务，只对当前局部轨迹进行更新。该类方法的主要问题是计算窗口与规划质量之间存在矛盾：预测范围过小可能导致局部最优，范围过大又会增加在线计算负担。

学习方法为动态规划提供了另一种思路。强化学习可以通过与环境交互学习状态和动作之间的映射，在复杂环境中快速给出局部动作\cite{Kober2013,westheider2023multiipp,ruckin2022adaptiveipp}。其优势在于可以把路径长度、能耗、碰撞风险以及任务收益等指标组合到奖励函数中，也能够表达一些难以准确建模的复杂策略。但工业巡检对安全性和可验证性要求较高，训练环境与现场环境之间的差异可能使策略性能下降。因此，在实际系统中，学习方法更适合与显式安全约束和传统规划结合，而不是完全替代约束优化。

总体来看，无人机巡检规划已经由点到点路径规划逐步发展到区域覆盖、目标视点规划、信息增益驱动规划和在线重规划。研究对象也从飞行几何约束扩展到观测质量和信息价值。但目前不少研究仍以预先确定的任务集合为前提，检测结果在任务执行过程中的反馈作用有限。对于自主巡检而言，更重要的问题是建立“观测结果—任务状态—下一观测行为”之间的闭环关系，使巡检轨迹能够随设备状态和信息不确定性变化而调整。

\subsection{多无人机任务分配与协同决策}

当巡检范围扩大、设备数量增加或任务具有时间要求时，单架无人机容易受到续航、载荷和任务时间的限制。多无人机系统能够并行执行任务，但与此同时产生了任务分配、任务排序、轨迹协同和避碰等问题。多机器人任务分配研究通常把任务收益、机器人能力以及任务之间的关系纳入统一模型\cite{Gerkey2004,Chakraa2023,ivic2023multiuavinspection,gilcastilla2024hetero,jeong2023makespan,obrien2023dynamictask}。对工业巡检而言，不同任务还可能对应不同传感器和观测条件，因此任务分配不仅是空间距离问题，也涉及平台能力与任务需求之间的匹配。

集中式任务分配利用中心节点掌握全部任务和无人机状态，通过整数规划或组合优化寻找整体较优方案。这种方式适用于规模较小、通信条件稳定的任务，但随着无人机和任务数量增加，组合复杂度以及中心节点的计算和通信压力都会增加\cite{Chakraa2023}。当任务不断变化时，每次重新求解整个优化问题又会带来额外延迟。因此，分布式任务分配逐渐成为多无人机研究的重要路线。

拍卖和一致性机制是典型的分布式方法。基于一致性的束算法（Consensus-Based Bundle Algorithm, CBBA）通过局部任务组合和一致性协商解决多个无人机竞争同一任务的问题，在不设置单一中央决策节点的条件下形成较一致的任务分配结果\cite{Choi2009}。这类方法便于把无人机位置、剩余电量、平台能力和任务收益写入局部竞价函数，也能够在任务集合发生变化时重新执行局部协商。因此，对于需要在线增加复查任务的工业巡检具有较好的适应性。

静态任务分配的局限在于通常假设任务集合基本稳定，而真实巡检中的任务可能来自实时感知。某处出现疑似裂纹后会新增复查任务，某个原定视点因临时障碍无法进入则需要重新安排任务，某架无人机剩余电量不足又可能触发任务转移。因此，动态任务分配需要能够响应任务到达、任务优先级变化以及无人机状态变化。近期综述已经将动态多无人机任务分配归纳为市场型、优化型和聚类型等主要方向，并把任务重分配、通信结构和动态事件处理作为重要研究问题\cite{Alqefari2025,obrien2023dynamictask,fei2022limitedcomm,Song2023MultiUAVMission}。

多无人机巡检中还存在异构资源问题。不同无人机可能具有不同的续航能力、载荷和传感器配置，适合执行的巡检任务也不相同。例如，红外检测任务需要具有热成像能力的平台，而三维测绘任务更适合搭载激光雷达的平台。近期研究开始采用聚类和分组机制减少候选无人机数量，从而降低任务分配和通信开销\cite{Dong2025,jacobsen2023swarm,gilcastilla2024hetero,jeong2023makespan}。这类方法比单纯按照距离分配任务更接近工业应用，但其效果依赖任务收益和平台能力评价是否合理。

任务分配与轨迹规划之间也存在明显耦合。距离某任务最近的无人机并不一定是最适合执行该任务的平台，因为加入新任务后可能造成任务序列改变、绕行增加或能源约束被突破。因此，一些研究已经开始同时处理任务分配和路径规划。针对输电线路巡检，Li等考虑三维山地环境中的多风场影响，将多无人机任务分配与路径规划联合起来进行优化\cite{Li2024Wind,ivic2023multiuavinspection}。针对风电机组协同巡检，Silano等通过信号时序逻辑描述任务和轨迹约束，在任务分配与轨迹生成之间建立统一关系，并进行了实际飞行验证\cite{Silano2024}。这些工作表明，多机协同已经由“先分任务、再规划路径”的串行结构向联合决策发展。

通信条件进一步限制了分布式协同。工业设施中的金属结构、建筑遮挡和地下空间可能导致通信链路不稳定，而连续传输图像和点云又会产生较高的数据负担。因此，多无人机系统需要考虑哪些状态必须共享、何时共享以及如何减少通信量。通信高效的CBBA分组方法表明，可以通过局部协商减少广播规模，同时保持任务分配的一致性\cite{Kim2020,fei2022limitedcomm,pan2024uhcrp}。但通信开销降低之后，系统获得的信息也可能变少，决策因此可能基于不完整或过期状态进行。由此，通信可靠性和任务分配质量之间存在需要共同处理的权衡。

多智能体强化学习也被用于动态任务分配和协同规划。相较于显式优化方法，它可以在训练阶段学习复杂的协同行为，并在在线阶段快速输出决策。但随着无人机和任务数量增加，联合动作空间会迅速扩大，训练稳定性和现实环境迁移问题也更加突出。对于工业巡检这种具有安全约束的任务，仅依靠经验策略难以形成充分的工程保证，因此较具应用前景的方案仍需要将学习策略与任务约束、能源限制和碰撞约束结合\cite{westheider2023multiipp,kulkarni2022teamed}。

综合上述研究，多无人机任务分配已经由集中式静态优化逐步发展到分布式协商、异构资源匹配、动态任务重分配以及任务和轨迹联合优化。现有研究较好地解决了“任务已经给定之后如何分配”的问题，但对“任务由感知结果动态产生”考虑仍不充分。对于复杂工业巡检，更合理的任务收益应该同时反映设备重要程度、异常风险、检测置信度、信息缺口以及无人机执行成本。只有当任务价值能够随设备状态变化而更新，多无人机协同才可能真正适应动态巡检过程。

\section{总结与展望}

前述研究从环境感知、设备异常认知和自主规划三个层面对复杂工业场景无人机巡检进行了梳理。现有研究已经形成较完整的技术基础：视觉、激光和惯性等传感器能够支持复杂环境中的定位和建图；监督式和无监督异常检测能够处理不同数据条件下的设备状态识别；覆盖规划、视点规划、在线重规划以及多无人机任务分配则能够支撑不同规模的巡检任务。问题在于，这些方法大多仍以模块为单位进行设计和评价，模块之间虽然可以通过接口连接，但异常检测和环境感知所产生的不确定性并没有充分进入后续规划和协同决策。

\subsection{现有研究存在的主要问题}

第一，环境感知的可靠性与后续巡检任务之间缺少统一联系。当前多模态感知方法主要使用定位误差、建图精度和实时性评价性能，而工业巡检真正关心的是当前观测是否足以支撑设备识别和缺陷判断。当视觉受到反光、弱纹理或遮挡影响时，单纯报告定位精度并不能说明当前视点是否仍然具有足够的巡检价值。因此，感知质量需要由传感器层面的指标进一步转化为面向任务的不确定性描述。

第二，异常检测结果尚未充分转化为新的巡检行为。第三章中的监督检测、无监督检测以及基础模型方法已经提高了未知异常的识别能力，但很多系统仍然将检测结果作为飞行结束后的分析结果。对于自主巡检，更有意义的是根据异常位置、置信度和风险等级决定是否复查、从何处复查以及采用何种传感器。这要求检测模块输出的不只是异常类别或分数，还应包含能够直接进入规划器的状态信息。

第三，现有巡检规划仍较多依赖预先确定的任务集合。覆盖路径和视点规划已经能够处理复杂设备表面，但任务通常在执行前就已经确定。实际现场的临时障碍、目标可见性变化和设备异常会使任务集合随时间变化，因此需要从固定航线转向任务在线生成和动态重规划。

第四，多无人机协同对感知结果的利用程度仍然有限。现有任务分配通常以空间距离、任务时间、能源和平台能力构造收益函数，而异常检测的置信度、设备历史状态和信息缺口较少被直接纳入任务价值。当新的复查任务产生后，还需要根据不同无人机的传感器配置和剩余资源进行动态分配，传统静态收益模型难以完全适应。

第五，真实工业场景验证仍不足。已有研究在公开数据集、仿真环境和实验室条件下取得了大量成果，但实际工业现场同时包含光照变化、遮挡、动态障碍、通信波动以及传感器局部退化等因素。系统级评价也往往被拆成定位、检测和规划三个相互独立的指标，难以回答一个完整巡检任务究竟是否更可靠、更高效。因此，面向真实任务的联合实验和统一评价仍需要加强。

\subsection{下一步研究方向}

结合前述问题，后续研究可以围绕“感知—认知—决策”一体化的自主巡检框架展开。研究重点不宜停留在单一检测器或路径规划器性能的继续提升，而应放在不同环节之间的信息传递和反馈机制上。

其一，建立面向任务的不确定性感知与多模态信息融合机制。针对视觉、激光、惯性和红外等模态在不同环境中的可靠性差异，可以根据当前任务需要对观测质量进行评估，并动态调整不同模态的作用。对于设备近距离缺陷检查，系统需要关注的不是单纯的定位精度，而是观测位置、视角和传感器状态是否能够支持稳定判断。这样可以把多模态融合从固定的数据组合转向面向任务的信息选择。

其二，建立由异常检测结果驱动的主动复查机制。当系统发现疑似异常时，根据异常置信度和设备风险自动生成复查任务，并利用设备几何和当前地图选择更合适的观测位置。复查后的新观测再用于更新异常判断，形成“发现异常—选择观测—重新判断”的闭环。这样可以使第三章中的异常检测结果直接进入第四章中的规划过程。

其三，研究信息价值驱动的动态巡检规划。对于不同设备和不同异常状态，可以根据风险、置信度和信息缺口定义动态任务价值，使无人机优先访问能够最大程度减少状态不确定性的区域。已经获得充分信息的位置可以减少重复观测，而信息不足或异常不明确的位置则提高优先级。这样，路径规划的目标可以从几何距离和覆盖率进一步转向单位飞行代价下的有效信息获取。

其四，建立面向动态任务生成的多无人机协同机制。当异常复查任务在线产生后，多无人机系统根据位置、剩余电量、载荷能力、通信条件以及任务风险重新分配任务。任务收益随设备状态变化而更新，任务分配和局部轨迹规划同步调整。对于大规模系统，可以采用分层结构，上层完成任务价值评估和资源分配，下层负责局部轨迹和安全避障，以降低实时计算压力。

其五，加强真实工业环境下的系统级验证。后续研究应进一步构建包含实际设备、不同传感器退化状态、动态障碍以及长期运行过程的测试环境，并将定位、检测、规划和任务完成效果放到同一评价体系中。除传统精度指标外，还应关注漏检率、异常复查成功率、任务完成时间、单位任务能耗、通信开销以及安全事件等指标，以评价方法的工程适用性。

总体来看，复杂工业场景无人机自主巡检正在由预设航线和单模块自动化向面向任务的闭环自主运行转变。现有研究已经分别解决了环境感知、异常检测和路径规划中的大量基础问题，但真正困难的部分逐渐转移到模块之间的协同：感知结果如何改变任务，异常判断如何改变观测，新的任务又如何在多无人机之间重新分配。后续研究若能够围绕这一链条建立统一的状态表示和决策机制，就有可能把已有的感知、检测与规划方法由相互衔接的功能模块进一步发展为面向复杂工业任务的自主巡检系统。

